\section{Results}
\label{sec:results}

We generated $N{=}100$ synthetic longitudinal patient profiles using SA and compared
them against the $K{=}23$ real patients across four validation levels. This sample
size provided approximately 4$\times$ amplification while remaining in the regime
where SA produces diverse, non-degenerate samples. As a baseline, we also generated
$N{=}100$ patients from a regularized multivariate normal (MVN) distribution fitted
to the same data.

\subsection{Marginal Plausibility}
\label{sec:results-l1}

We first assessed whether SA-generated patients reproduced per-feature summary
statistics across all three visits. The pooled median mean relative error
(MRE) across all 216 feature--visit entries was 1.2\% (95\% bootstrap CI:
1.0--1.6\%), with 193 of 216 entries (89\%) below 5\%
(Table~\ref{stab:mre-summary}), indicating that SA captured the central
tendencies of the real population. Per-visit MREs for representative
coagulation features are summarized in Table~\ref{tab:summary}, where
Factor~II and Factor~VIII fell below 2\%, antithrombin below 3\%, and
fibrinogen below 1\% across all three visits. We verified that synthetic
patients were not memorized copies of real patients using two complementary
metrics summarized in Table~\ref{stab:memorization}. The mean novelty score
($1 - \max_k \cos(\hat{\vxi}, \hat{\mathbf{m}}_k)$) was 0.50 at
$\beta = \beta^*$, indicating that generated samples were on average 50\%
away from their nearest stored pattern in angular distance. The median
nearest-neighbor distance from each synthetic patient to its closest real
patient was 6.14 standardized units, compared with 6.87 for the median
real-to-real nearest-neighbor distance, a ratio of 0.89 indicating that
synthetic patients were approximately as far from real patients as real
patients were from each other.

We next examined whether known physiological relationships and longitudinal
trajectories were preserved in the synthetic cohort. Pairwise scatter plots of
coagulation factor levels against thrombin generation parameters showed that
SA-generated patients occupied the same joint regions as real patients
(Fig.~\ref{fig:bio-correlations}). The expected inverse relationship between
antithrombin and TF-initiated peak thrombin was preserved, as were the positive
associations between Factor~VIII and peak, prothrombin and endogenous thrombin
potential (ETP), and the inverse relationship between antithrombin and ETP.
Per-visit means and standard deviations for six key coagulation features confirmed
that synthetic patients tracked the real longitudinal trajectories closely, with
overlapping confidence bands at all three visits
(Fig.~\ref{fig:pregnancy-progression}). Fibrinogen, Factor~VIII, and von Willebrand
factor showed the expected monotonic increases from baseline through the third
trimester in both cohorts, while Factor~II increased modestly and antithrombin
declined then partially recovered, consistent with the known hypercoagulable shift
of pregnancy. Together, these results indicated that SA captured not only univariate
distributions and pairwise dependencies but also the directionality and magnitude of
longitudinal change. The 72 features also included viscoelastic (ROTEM) and
fibrinolytic parameters, which were generated as part of the concatenated vector and
showed comparable MREs to the coagulation factors (Table~\ref{stab:mre-summary}), but are
not individually highlighted here because the mechanistic validation uses the BZ2012
thrombin generation model, which does not cover fibrinolysis or clot viscoelasticity.
\rone{These variables contributed through their loadings on the retained principal
components to the memory patterns that defined the energy landscape. Their fidelity
was evaluated statistically here, whereas their mechanistic fidelity was not evaluated.}

We also compared SA against two widely used deep generative methods for tabular data,
Conditional Tabular GAN (CTGAN) and Tabular Variational Autoencoder
(TVAE)~\citep{Xu2019CTGAN} (Table~\ref{stab:ctgan}). CTGAN failed completely,
producing median MREs of $\approx$19\%, an order of magnitude worse than SA, across
all epoch counts tested (300--3,000), consistent with known GAN mode collapse on small
datasets. TVAE achieved comparable marginal fidelity at high epoch counts (median MRE
1.8\% at 3,000 epochs), but operated on per-visit rows ($n{=}90$) rather than
concatenated longitudinal profiles ($n{=}23$), meaning it could not capture cross-visit
covariance structure and had no mechanism for conditional subgroup generation.

\rone{We next asked how much of this fidelity was attributable to the shared
principal-component representation rather than specifically to the stochastic-attention
generative formulation. We compared stochastic attention with four generators operating in
the identical $d{=}18$ PCA subspace: a two-component diagonal-covariance Gaussian mixture, a
kernel-density estimator, $k$-nearest-neighbor interpolation, and a Gaussian copula
(Supplementary Table~\ref{stab:pca-ablation}). The comparison indicated that the shared
representation accounted for much of the observed fidelity but did not guarantee
performance. The nearest-neighbor samples remained close to stored profiles: their median
novelty, defined as $1-\max_k\cos(\hat{\vxi},\hat{\mathbf{m}}_k)$, was $0.05$, and only
$5\%$ exceeded the operational $0.2$ cutoff, corresponding to cosine similarity below $0.8$
to every stored pattern. The kernel-density estimator had higher marginal error ($1.9\%$)
and lower pooled TF-only mechanistic cloud overlap ($0.82$): after pooling across patients,
visits, and TGA features, $82\%$ of the predicted-to-recorded TGA ratios for its synthetic
patients fell within the real-data 5th--95th percentile range. By contrast, the mixture
and copula were competitive with stochastic attention on marginal error ($1.1\%$ and
$1.3\%$, respectively, versus $1.2\%$) and mechanistic overlap ($0.92$ and $0.93$ versus
$0.90$). All $100$ profiles generated by stochastic attention and by the copula exceeded
the novelty cutoff, compared with $94\%$ of mixture profiles. Stochastic attention
nevertheless had the largest empirical cross-visit correlation error among the five methods
($0.64$, versus $0.42$ for the mixture and copula). Because this metric compared each
generated cohort with the same $23$ patients used to construct it, it measured in-sample
reproduction rather than population generalization. The separate known-ground-truth
benchmark addressed the latter question: stochastic attention recovered the population
covariance more accurately than the sample-covariance Gaussian generator in $35$ of $39$
low-rank Gaussian $n<p$ settings (Fig.~\ref{fig:sim-recovery}A). Thus, the ablation
attributed much of the marginal and mechanistic fidelity to the shared principal-component
representation. Within this comparison, no generator dominated across the evaluated metrics.}

\subsection{Cross-Visit Covariance Structure}
\label{sec:results-l2}

Having established that SA reproduces marginal distributions and pairwise
relationships, we next asked whether it preserves the higher-order joint structure
across visits. We compared the full $216 \times 216$ cross-visit correlation matrices for real,
SA-generated, and MVN-generated populations. The real data exhibited a characteristic
block structure with strong within-visit correlations along the diagonal and
structured off-diagonal blocks reflecting cross-visit dependencies; for example,
a patient's Visit~1 Factor~X level predicting their Visit~3 Factor~X level. We found
that SA preserved this block structure, including the off-diagonal cross-visit
correlations (Fig.~\ref{fig:cross-visit-corr}). The SA$-$Real residual matrix showed
small, unstructured deviations distributed throughout, whereas the MVN$-$Real residual
was notably smoother, reflecting the regularization-induced shrinkage of off-diagonal
correlations toward zero. This difference was most apparent in the off-diagonal blocks,
where MVN systematically underestimated cross-visit dependencies that SA retained.
We note that Pearson correlation captures linear associations; Spearman rank
correlations, used in the mechanistic validation, would additionally capture
monotonic nonlinear dependencies.

We examined the eigenvalue spectrum of the sample covariance matrix to understand the
structural origin of this difference (Fig.~\ref{sfig:eigenvalue-spectrum}). The real data had
rank~22, reflecting the $K{=}23$ patient constraint. SA and MVN make different
trade-offs in handling this rank deficiency: SA truncates via PCA (zeroing variance
beyond component~18), while MVN regularizes via Ledoit--Wolf shrinkage (providing a
full-rank estimate by inflating eigenvalues in the 194 null dimensions). The
consequence of MVN's approach is that it introduces variance in dimensions where the
data contain no signal, which manifested as increased dispersion in PCA projections.
PCA projections by visit confirmed this dispersion gap (\rtwo{Supplementary Fig.~\ref{sfig:pca-by-visit}}).
SA-generated patients occupied the same region of PCA space as real patients across
all three visits, while MVN-generated patients showed substantially greater scatter,
particularly at Visits~2 and~3.

\subsection{\rboth{Robustness to Sample Size, Dimensionality, and Nonlinearity}}

\rboth{We evaluated robustness with three complementary benchmarks
(Methods, Sec.~\ref{sec:robustness-benchmarks}). First, we simulated low-rank Gaussian
populations whose covariance was known before drawing any training data. For each combination
of training size ($n{=}8,15,23,40,80$), dimensionality ($p{=}30,120,216$), and latent rank
($r{=}3,5,10$), we constructed the stochastic-attention memory matrix and estimated a
sample-covariance Gaussian generator from the same training cohort, generated $100$ profiles
from each, and compared each generated
covariance with the known population covariance using relative Frobenius error. This tested
population recovery rather than reproduction of the training draw. Stochastic attention had
lower error in $35$ of the $39$ rank-deficient $n<p$ settings
(Fig.~\ref{fig:sim-recovery}A). At the cohort-matched setting ($n{=}23$, $p{=}216$), averaged
across the three latent ranks, the Gaussian generator's error was about $1.5$ times that of
stochastic attention. Across the grid, the advantage was concentrated in $n<p$ settings,
where the Gaussian sample covariance was rank deficient. Second, we tested degradation on
the clinical data by drawing ten random subsets at each of $n{=}20,15,10,$ and~$8$,
constructing a stochastic-attention memory matrix from each subset, and scoring each generated
cohort against the same full $23$-patient reference. As $n$ fell from $20$ to~$8$, mean
cross-visit correlation error rose
from $0.63$ to $1.11$ and mean pooled
marginal error rose from $2.0\%$ to $4.9\%$, while mechanistic cloud overlap remained near
$0.90$ (Fig.~\ref{fig:sim-recovery}B). Third, we tested the linear embedding against a known
nonlinear target: a noisy S-curve embedded in $p{=}30$ or~$120$ dimensions. We constructed
the stochastic-attention memory matrix and estimated the Gaussian generator from $n{=}23$
samples, then measured both distance from the known curve and covariance error as curvature
increased. The off-manifold error of stochastic attention doubled
relative to the straight-line case between curvatures $0.25$ and~$0.5$. Its samples remained
closer to the curve than the Gaussian samples, whereas the Gaussian generator recovered the
population covariance more accurately. Thus, stochastic attention outperformed this
sample-covariance Gaussian comparator on covariance recovery in most small, linear low-rank
settings, but the benchmarks also identified deterioration with fewer patients and a clear
limitation when the population structure was nonlinear.}

\subsection{\rtwo{Interpolation and Extrapolation}}

\rtwo{We next asked whether the sampler stayed within the real cohort or reached beyond it.
We constructed the convex hull of the $23$ real patients in the $d{=}18$
principal-component space and tested whether it contained each synthetic patient
(Supplementary Table~\ref{stab:hull-geometry}). This hull was the smallest convex region
enclosing the observed patients. Every synthetic patient lay outside it. In a separate
leave-one-out analysis, each real patient also lay outside the convex hull formed by the
remaining $22$ real patients. This indicated that every real patient formed an outer corner,
or vertex, of the sparse, high-dimensional hull. This binary inside-or-outside result did not
show whether synthetic patients lay just beyond the hull or much farther away. We therefore
used distance to the hull to measure the degree of extrapolation. Our hypothesis was that
unusually extrapolative generation would place synthetic patients farther from the hull than
real patients treated as unseen. For a like-for-like comparison, the distances for both groups
were calculated against hulls containing $22$ real patients. The leave-one-out calculation already
omitted the real patient being evaluated; for each synthetic patient, we omitted its nearest
real patient. Median distance to the resulting hulls was $11.0$ for synthetic patients and
$11.9$ for held-out real patients, and a descriptive Mann--Whitney test did not detect a
difference ($p{=}0.07$). Thus, the comparison provided no evidence that the synthetic
profiles lay farther beyond the observed cohort than held-out real patients. This result did
not establish equivalence, but it showed why binary hull membership alone was not an
informative measure of extrapolation in this sparse, high-dimensional setting.
}

\rtwo{Because geometric proximity did not establish biological validity, we next assessed
plausibility in two ways. First, $54$ of $100$ synthetic patients satisfied all implemented
biological constraints; the remaining $46$ had at least one negative estradiol or
progesterone concentration. The unconstrained concatenated multivariate-normal baseline had
a comparable failure rate, with negative hormone concentrations in $43$ of $100$ patients.
Thus, this limitation affected both unconstrained continuous generators rather than
stochastic attention alone. Second, we asked whether the generated coagulation-factor
combinations produced thrombin-generation behavior consistent with the real cohort under the
BZ2012 model. Estradiol and progesterone were not inputs to this model, and all nine generated
coagulation-factor inputs were positive; we therefore retained all $100$ synthetic patients
for this mechanistic check. For each thrombin-generation feature, $83$--$92\%$ of synthetic
patients had a model-predicted-to-recorded ratio within the 5th--95th percentile range
observed in the real cohort. Forty-one of $100$ patients met this range for every evaluated
feature at every visit, and $24$ of $100$ met both this criterion and all biological
constraints. While not included in the current implementation, future implementations could enforce non-negativity by construction through
positive-variable transformations or constrained decoding.}

\subsection{\rone{Membership-Inference Risk}}

\rone{Each real patient's $18$-dimensional principal-component profile was one column of the
memory matrix used during stochastic-attention sampling. We asked whether someone who already
possessed a candidate de-identified assay profile could use the synthetic cohort to infer
whether that profile was included in the matrix. This tested record inclusion, not personal
identity; the source cohort contained no direct identifiers or links to named individuals. As
an initial check, synthetic profiles were closer to their nearest real record than real
records were to one another (median distance $13.8$ versus $15.9$). We tested membership more
directly across $40$ random splits. For each split, the memory matrix contained $15$ patients,
the other $8$ were held out, and a new synthetic cohort was generated. A smaller distance from
a candidate record to that cohort provided stronger evidence of inclusion (Supplementary
Table~\ref{stab:membership}). Included records were closer than held-out records (median
$10.7$ versus $15.4$). The area under the receiver operating characteristic curve (AUC)
measured how often the test ranked an included record above a held-out record; $0.5$ indicated
chance and $1.0$ perfect discrimination. The mean AUC was $0.97$, with perfect discrimination
in $9$ of $40$ splits. Thus, the test usually inferred whether a candidate profile had been
included, but could not identify a person or reconstruct an unknown patient's measurements.
We next asked whether this signal came from the numerical sampler or the target distribution.
Metropolis-adjusted sampling left the result nearly unchanged (AUC $0.96$). Lowering the
inverse temperature across a $40$-fold range increased the distance from synthetic to stored
profiles, but it never reached the real-to-real baseline and cross-visit covariance fidelity
degraded. Exact draws from the closed-form mixture removed the sampling algorithm entirely
and reproduced the result (AUC $0.98$; Supplementary
Table~\ref{stab:sampling-ladder}, Fig.~\ref{sfig:sampling-ladder}). Thus, at $K{=}23$, the
membership signal came from the target distribution centered on the stored records, not from
how it was sampled. Synthetic generation therefore added no privacy guarantee beyond
de-identification of the source cohort.}

\subsection{Conditional Generation of Rare Subgroups}
\label{sec:results-l3}

The previous analyses used unconditioned SA, generating from the full cohort. A key
clinical application, however, is generating patients from specific subpopulations
that are too small to study independently. We therefore tested whether SA could
amplify rare clinical subpopulations while preserving their
condition-specific signatures. We defined three overlapping subgroups by pregnancy
outcome and comorbidity: Uncomplicated ($n{=}18$), PCOS ($n{=}3$, cross-cutting;
1~patient also in the PE group), and Developed~PE ($n{=}5$). The PCOS and
Developed~PE groups were too small for any independent statistical analysis. We used
SA's multiplicity-weighted sampling to generate condition-specific cohorts of 100
patients each by upweighting attention on the relevant stored patterns.

We compared condition-specific feature means between real and SA-generated patients
across eight coagulation features that exhibited between-condition variation
(Fig.~\ref{fig:conditioned-features}). SA-generated cohorts preserved the
condition-specific patterns. PCOS patients showed elevated Factor~VIII and vWF
relative to Healthy patients in both real and synthetic data, PE patients showed
elevated $\alpha_2$-AP and Factor~IX, and the between-condition rank ordering was
maintained across all eight features. We then performed a \rone{descriptive bootstrap
Mann--Whitney diagnostic}. For each feature and condition, we subsampled the
synthetic cohort to match the real sample size ($n_{\text{real}}$), applied a
Mann--Whitney U test, and repeated 1,000 times. We reported the fraction of
replicates where $p > 0.05$, \rone{meaning that this test did not detect a difference in
that replicate. This power-limited fraction was descriptive and did not establish
equivalence.} Across all 24 feature--condition pairs, 20 (83\%) achieved
$p > 0.05$ in $\geq$90\% of replicates, with a median
no-difference-detected fraction of 98.6\% (full per-pair results in
Table~\ref{stab:bootstrap-mw}). The four pairs that fell below 90\% were
high-variance features in the smallest subgroups (FIX and plasminogen in PCOS;
FIX and plasminogen in Developed~PE), where the real data exhibited large
inter-patient variability. \rone{No multiple-testing correction was applied because this
diagnostic was not used to make simultaneous inferential claims; the per-pair values and the
90\% threshold were descriptive reporting summaries.} PCA projections of the conditioned cohorts are provided in
Figs.~\ref{sfig:conditioned-pca}--\ref{sfig:conditioned-by-visit}. These results showed that SA's multiplicity-weighted
interpolation can amplify rare subgroups in a way that preserves their distinguishing
clinical features, useful for \rone{exploratory hypothesis generation and simulation-based
workflow testing, but not for increasing the effective clinical sample size}. MVN has no mechanism for this:
fitting a separate MVN to three PCOS patients across 216 features is mathematically
impossible (rank~2 covariance), and post-hoc subsetting of unconditional MVN samples
does not preserve subgroup-specific structure.

\rone{The conditioned cohorts reproduced the observed subgroup signatures, but generating
$100$ profiles did not increase the amount of independent subgroup information. The PCOS and
Developed~PE signals still came from only three and five real patients; after accounting for
the lower-weight background patterns, the corresponding participation ratios were $4.6$ and
$7.7$ (Table~\ref{stab:multiplicity}). Treating the $100$ synthetic profiles as independent
observations would therefore understate uncertainty. Conditioning instead provided a denser
set of profiles around the distributions suggested by those few patients, which could support
exploratory hypothesis generation, workflow testing, and simulations for planning future
studies. It did not add evidence about the conditions or replace recruitment of additional
patients.}

\subsection{Mechanistic Consistency}
\label{sec:results-l4}

The preceding validations assessed statistical properties of the synthetic data. We
next tested whether SA-generated patients produce biologically plausible outputs
when their coagulation factor levels are fed through
\rboth{a mechanistic model that knew nothing about the SA generation process, though it
was calibrated on the same cohort.}
We used the Hockin--Mann BZ2012 model~\citep{Hockin2002,BrummelZiedins2012}, a system
of 58 ordinary differential equations with 64 rate constants that simulates thrombin
generation from patient-specific coagulation factor inputs. We calibrated 5 of the
64 rate constants (prothrombinase, intrinsic and extrinsic tenase, protein~C
activation, and meizothrombin conversion; Table~\ref{stab:ode-calibration}) on Visit~1 real
patients at the population level and held the remaining 59 at published literature
values. We then ran the calibrated model on all 23 real patients and 100 synthetic
patients under both TF-only and TF+TM conditions (738 total simulations, 0 failures). For each patient, we
computed five ODE-predicted thrombin generation features (lagtime, peak, time-to-peak,
maximum rate, and ETP) and compared them against the corresponding values from the
patient record. In these comparisons (Fig.~\ref{fig:mechanistic-combined}, top row),
the horizontal axis is the TGA value from the patient record, which exists for both
real and synthetic patients, while the vertical axis is the value computed by the
BZ2012 model from that patient's factor inputs. Systematic deviations from the
$y{=}x$ line reflect ODE model bias (e.g., lagtime is underpredicted at
$\approx$0.77$\times$), not a failure of the synthetic data; the key observation is
that real and synthetic patients share the same pattern of model bias, occupying the
same cloud with the same systematic offset.

We formalized this observation by computing the ODE-predicted/dataset ratio for each
patient and comparing the resulting distributions between real and synthetic
populations (Fig.~\ref{fig:mechanistic-combined}, bottom row). These ratio
distributions capture how the ODE model processes each patient's factor levels. If
SA had generated patients with biologically implausible factor combinations, the
model would process them differently, producing a shifted or broadened ratio
distribution. Instead, the distributions overlapped substantially, with cloud
overlap (fraction of synthetic ratios within the 5th--95th percentile of real
ratios) ranging from \rone{0.83 for ETP to 0.92 for T.Peak} under TF-only conditions,
and two-sample Kolmogorov--Smirnov (KS) tests \rone{did not detect a difference between
the real and synthetic ratio distributions at this sample size} across all five TGA features
($D = 0.081$--$0.123$, all $p > 0.30$; full diagnostics in
Table~\ref{stab:mechanistic-validation}).
Under TF+TM conditions, the model systematically overcorrected the protein~C
anticoagulant pathway, producing peak and maximum rate predictions
approximately 0.5$\times$ measured values for both real and synthetic patients
(Fig.~\ref{sfig:mechanistic-tm-combined}, same layout as the main-text
comparison), yet this shared systematic bias was itself informative. Cloud overlap remained high (\rone{0.88--0.92} across the
five features), confirming that the model processed real and synthetic
patients identically even under conditions where the calibration was
imperfect.

The calibration also generalized across pregnancy timepoints. Although fitted
only on Visit~1 real patients, the per-visit median predicted-to-measured
ratios under TF-only conditions remained close to 1.0 at Visits~2 and~3
(Fig.~\ref{sfig:generalization}), confirming that the calibrated rate
constants captured a stable
population-level property of the coagulation system rather than overfitting to the
training visit. Spearman rank correlations between predicted and measured values were
moderate for ETP ($\rho \approx 0.6$--$0.8$ for real, $\rho \approx 0.6$--$0.65$
for synthetic; Fig.~\ref{sfig:rankcorr}) and weak for other features, consistent with
a 5-parameter population-level calibration that was not designed to resolve
inter-patient variability. The key finding was not patient-level prediction but
population-level plausibility. \rboth{A separately specified mechanistic model of the
coagulation cascade, applied identically to both groups and blind to the generator, placed
patients generated by stochastic attention within the same model-predicted envelope as real patients under both
experimental conditions. The model was calibrated on this same cohort, so this was a
same-cohort consistency check rather than independent validation.}

\subsection{Downstream Utility: Mechanistic Model Calibration}
\label{sec:results-utility}

The preceding results established that SA-generated patients are statistically and
mechanistically plausible. To test whether this plausibility translates to practical
utility, we calibrated a mechanistic model entirely on synthetic patients and
evaluated whether it could predict real patient outcomes. We calibrated the BZ2012 model
(TF-only condition, 5 rate constants) on synthetic V1 patients ($N{=}100$) using
the same bounded Nelder--Mead optimization applied to the real V1 calibration
($K{=}23$), with 11 random restarts to mitigate local-minimum sensitivity
(Table~\ref{stab:ode-calibration} for parameter details). We then evaluated both calibrations
on held-out real V2 and V3 patients that neither model had seen during
training. The synth-calibrated model achieved comparable or slightly better
performance across all five TGA features, with per-feature median relative
errors 2--10\% lower than the real-calibrated model (overall ratio
0.94$\times$; Fig.~\ref{fig:downstream-utility}; per-feature breakdown in
Table~\ref{stab:downstream-utility}). This improvement likely
reflects the smoother loss landscape afforded by $N{=}100$ synthetic training patients compared with
$K{=}23$ real ones, reducing overfitting during optimization. The two calibrations
found different parameter values but produced highly correlated predictions across all
five TGA features, confirming that the synthetic data captured the same underlying
coagulation structure. We note that the synthetic training data is one step removed
from the real data (real $\to$ SA $\to$ synthetic $\to$ calibration), so the
synth-calibrated model is not fully independent of the real data; rather, it
demonstrates that SA's representation preserves sufficient biological structure for a
mechanistic model to learn generalizable parameters.

\rone{The synthetic-calibrated model's median prediction errors were 2--10\% lower than those
of the real-calibrated model, but this was not interpreted as evidence that synthetic data
were superior. All information in the synthetic
cohort derived from the same $23$ patients, and calibrating five
rate constants against $100$ samples drawn from that distribution likely smoothed the
objective and diluted the influence of any single patient. Any generator producing a denser
sample of the same distribution could have had the same variance-reduction effect. The result
therefore showed that the synthetic cohort was usable for calibration, not that it contained
new biological information.} \rtwo{Together with the conditioned-cohort analyses, this
defined the method's practical value: it supported mechanistic-model calibration without
using the real cohort directly in that calibration, comparisons of subgroups containing only
three or five patients, and factor-level hypotheses for testing in a larger study. These uses
did not replace additional patients or support clinical deployment, which would require
prospective validation against real clinical decisions.}
